\section{Grupos y Homomorfismos}

\newcommand{\comentario}[1]{\textcolor{gray}{\(\rightarrow \) #1}}
\subsection*{Permutaciones}
\begin{definition}
    Sean \(S_X:= \{\alpha: X \to X \mid \alpha \text{ es biyección}\}\) el conjunto de todas las permutaciones de \(X\) y \(*: S_X \times S_X \to S_X\) tal que \(\alpha * \beta := \alpha \circ \beta\).  En la práctica,   
    \[ X := \{1, 2, \ldots, n\}\text{ \  e \ } \alpha =  \begin{pmatrix}
        1 & 2 & \cdots & n \\ 
        \alpha(1) & \alpha(2) & \cdots & \alpha(n)
    \end{pmatrix}.\]
\end{definition}
\vspace{-0.7cm}
\begin{definition}
    \(\alpha  \in S_n\) es un $r$-\emph{ciclo} si fija \(n-r\) posiciones e \(\alpha :  i_1 \mapsto i_2 \mapsto \cdots \mapsto i_r \mapsto i_1\), escribimos \(\alpha = (i_1 \ \  i_2 \ \  \ldots \ \  i_r)\). Para los casos \(r=1, 2\) tenemos
    \[\star \ \  (i_j) =: \id_X  \ \leftarrow \text{ \emph{Identidad} } \hspace{2cm}\star \ \ (i_1 \ \ i_2) \ \leftarrow \text{ \emph{Transposición}}\]
\end{definition}

\ejemplo{
    \begin{example}
        \(S_3 \ni \alpha \beta = \begin{pmatrix}
            1 & 2 & 3 \\ 
            3 & 2 & 1  
        \end{pmatrix} \begin{pmatrix}
            1 & 2 & 3 \\ 
            2 & 3 & 1
        \end{pmatrix} = 
        (1 \ \ 3 )( 1 \ \ 2 \ \ 3) 
        = (3)(1 \ \ 2) \). 
    \end{example}
}

\begin{definition}
    Dos ciclos \(\alpha, \beta \in S_X\) son \emph{disjuntos} cuando afectan diferentes elementos. En otras palabras \(\{x \in X : \alpha (x) \neq x  \text{ \ y \ } \beta(x)\neq x\}= \varnothing\). %A family \((\alpha_j)\) is disjoint if each pair of them is disjoint.   
\end{definition}

\begin{theorem}
    \textbf{1.1.} Toda \(\alpha \in S_n\) es ciclo o producto de ciclos disjuntos. %is either a cycle or a product of disjoint cycles. 
\end{theorem}

\demo{Por inducción en \(k\) (elementos permutados). Sea \(i_1\) elemento permutado por \(\alpha\), planteé el ciclo \(\sigma = (i_1 \ \ i_2 \ \  \ldots \ \ i_r)\). Sigue que \(\alpha = \sigma \alpha'\), donde \(\sigma \) e \(\alpha'\) son disjuntos e \(\alpha'\) permuta \(k - r < k\) elementos. }

\begin{definition}
    Una \emph{factorización completa} de \(\alpha\in S_n\) es una expresión de \(\alpha\) como producto de ciclos disjuntos que incluye \(1\)-ciclos (aunque no se escriban). 
\end{definition}

\begin{theorem}
    \textbf{1.2.} Sea \(\alpha \in S_n\). Factorizaciones completas \(\alpha = \beta_1\beta_2 \ldots\beta_k\) son únicas a menos de reordenación de los factores. %\comentario{\textcolor{rojoscuro}{Ex. 6, 8 and 14}}
\end{theorem}

\demo{Sea \(\alpha = \gamma_1 \gamma_2 \ldots \gamma_s\) una segunda factorización completa. Si \(\beta_t\) permuta \(i_1\), entonces \(\exists \gamma_j \) que permuta \(i_1\), sea \(j=s\) (\textcolor{rojoscuro}{Ex. 1.8}). Ahora, \(\forall k \in \N,\ \beta_t^k(i_1) = \alpha^k(i_1) = \gamma_s^k(i_1)\) (\textcolor{rojoscuro}{Ex. 1.14}), luego \(\beta_t = \gamma_s\) e \(\beta_1 \ldots \beta_{s-1} = \gamma_1\ldots \gamma_{t-1}\) (\textcolor{rojoscuro}{Ex. 1.6}). El resultado sigue por inducción en \(\max\{s,t\}\). }

\begin{theorem}
    \textbf{1.3.} Toda \(\alpha \in S_n\) es producto de transposiciones. \comentario{Por el \(\blacksquare\) \hyperlink{thm11}{1.1.} basta mostrar para ciclos, observese que \((1 \ \ 2 \ \ \ldots \ \ r) = (1 \ \ r) (1 \ \ r-1)\cdots(1 \ \ 2)\)}
\end{theorem}



\newcommand{\sgn}{\operatorname{sgn}}
\begin{definition}
    Dada \(\alpha = \beta_1 \ldots\beta_t  \in S_n\) (factorización completa), \(\sgn(\alpha):= (-1)^{n-t}\). Tambien, \(\alpha \) es \emph{par} (o \emph{impar}) si es producto par (o impar) de transposiciones.
%    \[\star \ \ \sgn(\alpha) = 1 \ \leftarrow \ \alpha \text{ \emph{par}} \hspace{2cm} \star \ \ \sgn(\alpha) = -1 \ \leftarrow \ \alpha \text{ \emph{impar}}\]
\end{definition}

%\begin{note}
%    La paridad de \(\alpha \) se corresponde con la cantidad (par u impar) de factores en su descomposición en transposiciones.       
%\end{note}

\begin{theorem}
    \textbf{1.7.} Para todas \(\alpha, \beta \in S_n, \ \sgn(\alpha\beta)= \sgn(\alpha) \sgn(\beta)\). Más aún, \(\alpha\) es par (o impar)\ \(\leftrightharpoons\) \ \(\sgn(\alpha) = 1\) (o \(-1\)). \comentario{No admite resumen, \cite[pag. 9]{stein}}
\end{theorem}

\subsection*{Grupos}

\begin{definition}
    Un \emph{grupo} \((G, *)\) es un conjunto no vacío \(G\), junto con una operación binaria asociativa \(* : G \times G \to G\) tal que: 
    \vspace{-0.2cm}
    \begin{enumerate}[left = 1.33cm]
        %\item[G1.] \(\forall a,b,c \in G, \ (a* b)*c = a * (b*c)\). 
        \item[G1.] \(\exists e \in G : \forall a \in G, \ a*e = e *a = a\). \comentario{Elemento identidad}
        \item[G2.] \(\forall a \in G, \ \exists b \in G  :  \ a*b = b *a = e\). \comentario{Existencia de inversos}
        \item[+ G3.] \(\forall a,b \in G, \ a*b = b *a  \ \leftarrow\) \ \emph{Grupo abeliano}.
    \end{enumerate}
\end{definition}

\newcommand{\F}{\mathbb{F}}
\newcommand{\Z}{\mathbb{Z}}
\newcommand{\GL}{\operatorname{GL}}
\newcommand{\SL}{\operatorname{SL}}
\ejemplo{
   \textcolor{verdeoscuro}{\underline{\textbf{Examples}}} \ \  {\scriptsize\(\bullet\)} \ \((S_n, \circ)\)  \emph{grupo simétrico}. \hfill \(\mid\) \hfill  {\scriptsize\(\bullet\)} \ \((\F, +)\), cuerpo con la suma. \vspace{0.3cm} \\ 
   {\scriptsize\(\bullet\)} \ \((\Z_n,+) := \{[a] \in \nicefrac{\Z}{\sim} :  a\sim b \ \leftrightharpoons \ a \equiv b \ (\operatorname{mod} n)\}\).  \hfill \(\mid\) \hfill {\scriptsize\(\bullet\)}  \ \((\GL_n(\F), \cdot)\).
}

\begin{theorem}
    \textbf{1.10.} Si \(G\) es un grupo, identidad e inversos son únicos. \comentario{Suponga que existe otro que satisface G\(1\) (o G\(2\)) y proceda como siempre }
\end{theorem}

\begin{corollary}
    \textbf{1.11.} Denotamos por \(a^{-1}\) al inverso de \(a\) en \(G\), y a la identidad como \(\id_{_G}\). Sigue inmediatamente que \((a^{-1})^{-1} = a\). 
\end{corollary}

\begin{definition}
    Dados \(a \in G, \ n \in \N\), definimos \(a^n := \overbrace{a * a * \cdots * a}^\text{$n$ \text{ vezes}}\) e \(a^{-n} := (a^{-1})^n\). Por convención \(a^0 := \id\), y llamamos a todas estas expresiones de \emph{potencias} de \(a\).
\end{definition}

\subsection*{Homomorfismos}
%
%If \(G = \{a_1 , a_2, \ldots, a_n\}\), we can write \(a_i * a_j = a_ia_j\) and draw a \emph{multiplication table}:  
%\begin{table*}[!h]
%\centering    
%\begin{tabular}{c | c c c c}
%    \(*\) & \(a_1\) & \(a_2\) & \(\cdots\)& \(a_n\)  \\ \hline
%    \(a_1\) & \(a_1a_1\) & \(a_1a_2\) & \(\cdots\)& \(a_1a_n\)\\ 
%%    \(a_2\) & \(a_2a_1\) & \(a_2a_2\) & \(\cdots\)& \(a_2a_n\)\\ 
%    \(\vdots\) & \(\vdots\) & \(\vdots\) & \(\ddots\)& \(\vdots\)\\ 
%    \(a_n\) & \(a_na_1\) & \(a_na_2\) & \(\cdots\)& \(a_na_n\)\\ 
%\end{tabular}
%\end{table*}
\begin{definition}
    Sean \((G, *)\) e \((H,\circ )\) grupos. Una función \(f: G \to H\) es \emph{homomorfismo} si \(\forall a,b \in G\), \(f(a*b) = f(a) \circ f(b)\). Si \(f\) fuera bijectiva, \(f:G\cong H\) (\emph{isomorfismo}).  
\end{definition}

\begin{theorem}
    \textbf{1.13.} Sea \(f : (G,*) \to (H, \circ)\) homomorfismo. Entonces,  
    \vspace{-0.3cm}
    \begin{enumerate}[label = \roman*., left = 0.3cm]
        \item \(f(\id_{_G}) = \id_{_H}\). \hfill  ii. \ \(f(a^{-1}) = f(a)^{-1}\). \hfill  iii. \ \(f(a^n) = f(a)^n\). 
    \end{enumerate} 
\end{theorem}

\demo{Para (i) e (ii) calcula \(f\) en \(\id_{_G} = \id_{_G} * \id_{_G}\) e \(a*a^{-1} = \id_{_G} = a^{-1} * a\). La afirmación (iii) sigue por inducción en \(n\) e (ii).}
